\section{Introduction}\label{sec:introduction}

Let $(X,\mathcal F,\mu)$ be a probability space and let $T:X\to X$
be an invertible measure-preserving transformation with measurable
inverse. Given bounded measurable functions $f_1,f_2,f_3$ and distinct
nonzero integers $\lambda_1,\lambda_2,\lambda_3$, consider
\begin{equation}\label{eq:averages}
 A_N(x)=\frac1N\sum_{n=1}^N
       \prod_{j=1}^3 f_j(T^{\lambda_jn}x).
\end{equation}
The issue is convergence at almost every point, rather than only in
norm. We assume that $T$ is mixing, in the sense that
\begin{equation}\label{eq:mixing}
 \mu(A\cap T^{-n}B)\longrightarrow\mu(A)\mu(B)
 \quad\text{as }\abs n\longrightarrow\infty
 \quad(A,B\in\mathcal F).
\end{equation}

\begin{theorem}\label{thm:main}
For every system satisfying~\eqref{eq:mixing}, every triple of pairwise
distinct nonzero integers $\lambda_1,\lambda_2,\lambda_3$, and every
bounded measurable $f_1,f_2,f_3:X\to\C$, the averages
in~\eqref{eq:averages} satisfy
\[
 A_N(x)\longrightarrow\prod_{j=1}^3\int_X f_j\,d\mu
 \quad\text{for almost every }x,
\]
as $N$ tends to infinity through all positive integers. The exceptional
null set may depend on the slopes and the functions.
\end{theorem}

The probability space need not be standard or countably generated.
Both positive and negative slopes are allowed. The proof depends on
the analytic results of the companion manuscripts~\cite{H,V}; their
role and the extension required here are described below.

\subsection{Context and predecessors}
Birkhoff's pointwise ergodic theorem~\cite{Birkhoff} treats one
function. Products along several iterates arise naturally in
Furstenberg's ergodic proof of the multiple recurrence theorem and its
connection with arithmetic progressions~\cite{Furstenberg}. Multiple
recurrence, norm convergence, and pointwise convergence are different
questions. In particular, convergence in $L^2$ does not by itself
establish convergence at almost every point along the full sequence.

For weakly mixing transformations, the product-of-integrals limit in
$L^2$ is classical. Bergelson's
weakly mixing PET theorem~\cite[Theorem~1.2 and p.~340]{Bergelson}
includes distinct nonzero integer slopes of either sign.
For general systems, Host and Kra~\cite{HostKra} and
Ziegler~\cite{Ziegler} developed the norm-convergence theory using
characteristic factors with nilpotent structure. These results explain
the limiting structure without resolving the corresponding pointwise
question.

Bourgain's double recurrence theorem~\cite{Bourgain} establishes
almost-sure convergence for two bounded functions evaluated at two
powers of one transformation. Pointwise theorems for longer products
also have substantial precedents under additional assumptions on the
system. For example, Derrien and Lesigne~\cite[Theorem~4.2]{DerrienLesigne}
prove polynomial pointwise theorems for $K$-systems and exact
endomorphisms; their invertible case includes the linear configurations
considered here. Other structural hypotheses yielding pointwise
convergence for all finite progression lengths include singular spectrum
on the Pinsker factor of a weakly mixing system, by
Assani~\cite{Assani}, and the pairwise-independent-joining property for
weakly mixing systems, by Gutman, Huang, Shao, and Ye~\cite{GHSY}.
The significance of Theorem~\ref{thm:main} is that its
dynamical assumption is mixing alone, while the three slopes are
arbitrary distinct nonzero integers.

The harmonic-analytic method separates convergence from identification
of the limit. A uniform oscillation estimate gives pointwise
convergence; mixing identifies that limit by an $L^2$ argument.
Oscillation methods in ergodic theory are developed in
\cite{Bourgain,JKRW}, and the passage from translation estimates to
orbit estimates follows the transference principle of
Calder\'on~\cite{Calderon}. We give that passage directly, so it imposes
no additional structure on the probability space.

\subsection{Analytic input and the coefficient extension}
The manuscript~\cite{H} proves an $L^3$ bound for the trilinear Hilbert
transform through local four-linear forms. Its proof includes chart
detection, localization by quadratic data, comparison on short depth
blocks, and a separated-frequency counting estimate. The quantitative
inverse theorem of Leng, Sah, and Sawhney~\cite{LSS} enters its chart
detection argument. Arithmetic regularity and inverse-theorem methods
also occur in Tao's sublogarithmic cancellation theorem for truncated
multilinear Hilbert transforms~\cite{Tao}; that cancellation theorem is
distinct from the uniform local estimate used here.

The companion~\cite{V} adapts the local estimates to inputs that vary
with depth and deduces pointwise convergence for the slopes $1,2,3$.
The local forms in both companions use the four coordinates
$x,x+t,x+2t,x+3t$. A power of $T$ and a permutation of the functions
do not turn an arbitrary triple into this configuration. Thus their
stated theorems cannot simply be invoked for the slopes in
Theorem~\ref{thm:main}.

The new work is the extension of their local proofs to
$x+b_jt$, $0\le j\le3$, for fixed distinct integers with $b_0=0$.
Three features are particularly useful beyond this application.
First, an integer left inverse of the \emph{whole} primitive
configuration replaces a unimodular pair of coordinates; this removes
extra lattice or torus configurations in discretization.
Second, a Vandermonde relation provides the quadratic cancellation and
the Fourier randomization for every set of four distinct slopes.
Third, a height base divisible by all pair differences preserves the
divisibility required in the separated-frequency argument. The last
step must preserve a bound polynomial in the dimension: an exponential
loss there would not fit the proof's parameter choices.

These are extensions of the proofs in~\cite{H,V}, not new proofs of
their coefficient-independent projection and row estimates. We state
the local conclusion completely, prove the coefficient-sensitive
identities and estimates, and specify the results imported from the
companions at their points of use. In particular, the technical part
is a proof by reference to those manuscripts, whereas the
ergodic deduction is presented in full.

\subsection{Proof structure}
Section~\ref{sec:local} defines the admissible kernels and states the
local estimate for inputs with controlled variation on blocks of
depths. Section~\ref{sec:transfer} proves Theorem~\ref{thm:main} from
that estimate: smooth scale differences are realized by cancellative
kernels, then transferred to integer sequences and finite orbits.
The resulting oscillation bound gives convergence for smooth profiles;
moment-corrected approximations to interval indicators recover all
integer lengths.

The remaining sections prove the local estimate. Section~\ref{sec:analytic}
establishes the coefficient modifications in detection and
localization. Section~\ref{sec:arithmetic} proves the short-block,
height, and counting modifications. Section~\ref{sec:completion}
closes the parameter argument, first for fixed inputs and then for
depth-dependent inputs. All analytic constants may depend on the fixed
slopes; no bound is required uniformly over triples of integers.
